\section{More Discussion on the Scaling Law Shapes}
\label{sec:scaling-law-discussion}
\label{sec:scurve-mechanistic-reading}

\underline{\textit{A mechanistic reading selects the log-sigmoid.}}
Because the candidate S-curves fit almost equally well
(Table~\ref{tab:scurve-family-comparison}; the log-sigmoid/log-probit degeneracy
has been known since Berkson~\citep{berkson1944application}), the choice cannot
rest on fit; we make it on mechanism. Written as a growth law for the normalized
score $y=S/S_{\max}$, each candidate's rate makes a different statement about
what drives progress and what limits it.
\begin{itemize}
  \item \textbf{Log-sigmoid:} $\frac{dy}{d\tau}=\beta\,y(1-y)$ in log time
  $\tau=\ln t$. The detailed derivation is given in
  Section~\ref{sec:why-log-sigmoid} and Appendix~\ref{sec:theory}; here the key
  mechanistic reading is that \(y\) is the unlocked attainable score mass and
  \(1-y\) is the remaining locked score mass. If the aggregate influence crossing
  the unlocked--locked frontier is approximately proportional to the product of
  these masses, then progress follows \(y(1-y)\). The log-time coordinate comes
  from the self-similar search geometry of the task graph. This reading is
  consistent with two empirical checks: Section~\ref{sec:stateful-stateless-analysis}
  shows that continuous stateful runs outperform equal-budget repeated sampling,
  so progress is not explained by independent attempts alone; and
  Section~\ref{sec:case-study} shows a finite task advancing through sparse but
  cumulative breakthroughs, where an early working pipeline makes later repairs
  searchable. Inflection occurs at \(y=0.5\) (symmetric).

  \item \textbf{Log-Gompertz:} $\frac{dy}{d\tau}=c\,y\ln(1/y)$. The
  $\ln(1/y)$ term is the signature of an engine that winds down
  multiplicatively as the system matures (e.g., tumour growth, where the
  proliferating fraction shrinks). It is front-loaded, with inflection at
  $y=1/e\approx0.37$: the relative growth rate diverges as $y\to0$, so a tiny
  system grows explosively because the limitation it will eventually hit does
  not yet exist. Experience acquisition is the opposite---its early phase is
  \emph{slow} because a foothold must first be bootstrapped, not fast for lack
  of a brake.

  \item \textbf{Log-Probit:} $\frac{dy}{d\tau}=\frac{1}{\sigma}\,
  \varphi\!\big(\Phi^{-1}(y)\big)$, a sweep across a \emph{log-normal}
  distribution of difficulties---difficulties formed as products of many
  \emph{independent} factors (a multiplicative central-limit argument).
  Experience acquisition is path-dependent rather than a product of independent
  factors, so this microfoundation does not apply, even though probit and the
  logistic are empirically near-indistinguishable.

  \item \textbf{Weibull CDF:} $y(t)=1-\exp\{-(t/\lambda)^\beta\}$, equivalently
  $\frac{dy}{dt}=h(t)(1-y)$ with
  $h(t)=\frac{\beta}{\lambda}(t/\lambda)^{\beta-1}$. Its hazard is a function
  of raw elapsed time, while accumulated progress enters only through the
  survival term $(1-y)$. This makes it the natural first-passage or repeated
  sampling baseline. For a single 0--1 reward task with independent attempts of
  success probability $p$,
  $P_{\rm pass}(k)=1-(1-p)^k=1-\exp[-k(-\ln(1-p))]$, which is an exponential
  CDF, the $\beta=1$ Weibull case (approximately $1-\exp(-pk)$ for small
  $p$). Such a curve improves because repeated independent attempts shrink the
  remaining failure mass, not because state carried across attempts makes later
  attempts more effective. Section~\ref{sec:stateful-stateless-analysis} shows
  that stateful learning beats this repeated-sampling mechanism under the same
  total budget, and the Weibull CDF also fits worse than the log-sigmoid in
  Table~\ref{tab:scurve-family-comparison}. Individual improvements may have a
  first-passage flavor, but the macroscopic best-so-far trajectory accumulates
  through dependent, stateful steps rather than through a raw-time hazard with no
  accumulated-progress factor.

  \item \textbf{Log-linear:} $\frac{dy}{d\tau}=b$ (constant). With no
  saturating term it grows without bound and cannot level off, contradicting the
  many tasks that reach a ceiling within the budget; it attains by far the worst
  fit.
\end{itemize}
We therefore read the log-sigmoid not merely as the best-fitting S-curve but as
the curve whose \(y(1-y)\) rate law matches the frontier-expansion interpretation:
unlocked score mass supplies reusable capability, while locked score mass bounds
the remaining opportunity for improvement.
This is a mechanistic preference, not an empirical exclusion---the data cannot
separate the symmetric families. It is falsifiable through the inflection: a
symmetric peak near $y=0.5$ is consistent with the logistic (and probit), a
front-loaded peak near $0.37$ would favor Gompertz, and a back-loaded peak near
$0.63$ would favor Weibull.
